\chapter{D\'etection et traitement des valeurs aberrantes}
\label{ch:outliers}

\begin{quote}
\emph{<< Une valeur aberrante n'est pas du bruit tant que vous n'avez pas prouv\'e que c'est du bruit. >>}
\end{quote}

% --- hook ---
Un patient pesant 800\,kg est une erreur de saisie. Une maison à 10 millions de dollars dans un quartier où la médiane est 500\,000 est peut-être réelle. Une valeur aberrante n'est pas du bruit tant que vous n'avez pas prouvé que c'en est. La supprimer sans investiguer corrompt votre analyse aussi sûrement que la garder. Ce chapitre couvre les outils de détection (z-score, IQR, Mahalanobis, Isolation Forest), mais surtout la discipline qui doit les accompagner.

% ============================================================
\section{Qu'est-ce qu'une valeur aberrante ?}

Une valeur aberrante (outlier) est une observation qui s'\'ecarte substantiellement du motif \'etabli par la majorit\'e des donn\'ees. Les outliers peuvent \^etre :

\begin{itemize}
  \item \textbf{Erreurs de saisie} : un poids de patient de 800\,kg (au lieu de 80).
  \item \textbf{Artefacts de mesure} : un capteur renvoyant $-999$ pendant la calibration.
  \item \textbf{Vraies valeurs extr\^emes} : une maison vendue 10 millions de dollars dans une banlieue o\`u la m\'ediane est 500\,000\$.
\end{itemize}

\begin{warningbox}
Ne supprimez jamais d'outliers automatiquement. Chaque outlier doit \^etre investigu\'e. Supprimer de vraies valeurs extr\^emes parce qu'elles << ont l'air fausses >> corrompt votre analyse. Supprimer des erreurs de donn\'ees parce qu'elles << ont l'air inhabituelles >> est correct et n\'ecessaire.
\end{warningbox}

% ============================================================
\section{D\'etection visuelle}

\begin{minted}{python}
import pandas as pd
import matplotlib.pyplot as plt
import seaborn as sns

# Le dataset FIFA 21 contient beaucoup de valeurs extr\^emes et d\'esordonn\'ees
df = pd.read_csv("fifa21_raw.csv")

# Les box plots r\'ev\`elent les outliers imm\'ediatement
fig, axes = plt.subplots(1, 3, figsize=(14, 5))

sns.boxplot(y=df["Age"], ax=axes[0])
axes[0].set_title("Age distribution")

sns.boxplot(y=df["Overall"], ax=axes[1])
axes[1].set_title("Overall rating")

sns.boxplot(y=df["Value(in Euro)"], ax=axes[2])
axes[2].set_title("Market value (Euro)")

plt.tight_layout()
plt.savefig("outlier_boxplots.png", dpi=150)
plt.show()
\end{minted}

\begin{minted}{python}
# Scatter plot : outliers dans l'espace 2D
fig, ax = plt.subplots(figsize=(8, 6))
ax.scatter(df["Age"], df["Overall"], alpha=0.3, s=10)
ax.set_xlabel("Age")
ax.set_ylabel("Overall Rating")
ax.set_title("Age vs Overall --- look for isolated points")
plt.tight_layout()
plt.show()
\end{minted}

% ============================================================
\section{M\'ethodes statistiques}

\subsection{M\'ethode du z-score}

\begin{minted}{python}
import numpy as np
from scipy import stats

# Z-score : nombre d'\'ecarts-types par rapport \`a la moyenne
z_scores = np.abs(stats.zscore(df["Overall"].dropna()))
threshold = 3
outliers_z = (z_scores > threshold)
print(f"Outliers by z-score (|z| > {threshold}): {outliers_z.sum()}")
\end{minted}

\begin{minted}{python}
# Appliquer le z-score \`a plusieurs colonnes
def detect_zscore_outliers(df, columns, threshold=3):
    """Return a boolean mask of rows with any z-score outlier."""
    mask = pd.Series(False, index=df.index)
    for col in columns:
        z = np.abs(stats.zscore(df[col].dropna()))
        col_mask = pd.Series(False, index=df.index)
        col_mask.iloc[:len(z)] = z > threshold
        mask = mask | col_mask
    return mask

num_cols = ["Age", "Overall", "Potential"]
outlier_mask = detect_zscore_outliers(df, num_cols)
print(f"Rows with at least one outlier: {outlier_mask.sum()}")
\end{minted}

\begin{datatip}
La m\'ethode du z-score suppose une distribution \`a peu pr\`es normale. Pour des donn\'ees asym\'etriques (revenus, prix immobiliers, valeurs de march\'e de joueurs), utilisez la m\'ethode IQR ou transformez en log d'abord.
\end{datatip}

\subsection{M\'ethode IQR}

\begin{minted}{python}
def iqr_bounds(series):
    """Return lower and upper IQR bounds."""
    Q1 = series.quantile(0.25)
    Q3 = series.quantile(0.75)
    IQR = Q3 - Q1
    lower = Q1 - 1.5 * IQR
    upper = Q3 + 1.5 * IQR
    return lower, upper

lower, upper = iqr_bounds(df["Overall"])
outliers_iqr = df[(df["Overall"] < lower) | (df["Overall"] > upper)]
print(f"IQR bounds: [{lower:.1f}, {upper:.1f}]")
print(f"Outliers: {len(outliers_iqr)}")
\end{minted}

\begin{minted}{python}
# Rapport IQR pour toutes les colonnes num\'eriques
def iqr_outlier_report(df, columns):
    """Generate an IQR outlier report."""
    rows = []
    for col in columns:
        lo, hi = iqr_bounds(df[col].dropna())
        n_out = ((df[col] < lo) | (df[col] > hi)).sum()
        rows.append({"column": col, "lower": lo, "upper": hi,
                      "n_outliers": n_out,
                      "pct": round(n_out / len(df) * 100, 1)})
    return pd.DataFrame(rows)

report = iqr_outlier_report(df, ["Age", "Overall", "Potential"])
print(report)
\end{minted}

% ============================================================
\section{D\'etection multivari\'ee}

\subsection{Distance de Mahalanobis}

\begin{minted}{python}
from scipy.spatial.distance import mahalanobis
from numpy.linalg import inv

def mahalanobis_outliers(df, columns, threshold=3):
    """Detect multivariate outliers using Mahalanobis distance."""
    data = df[columns].dropna()
    mean = data.mean().values
    cov_matrix = data.cov().values
    cov_inv = inv(cov_matrix)

    distances = data.apply(
        lambda row: mahalanobis(row.values, mean, cov_inv), axis=1
    )
    return distances, distances > threshold

cols = ["Age", "Overall", "Potential"]
distances, mask = mahalanobis_outliers(df, cols)
print(f"Multivariate outliers: {mask.sum()}")
\end{minted}

\begin{preprocesstip}
La distance de Mahalanobis d\'etecte des points qui ne sont aberrants sur aucune dimension individuelle mais qui sont inhabituels dans la \emph{combinaison} de dimensions. Un joueur de 40 ans est normal. Un joueur not\'e 95 est normal. Un joueur de 40 ans not\'e 95 est aberrant dans l'espace joint.
\end{preprocesstip}

\subsection{Isolation Forest}

\begin{minted}{python}
from sklearn.ensemble import IsolationForest

# Isolation Forest : d\'etection d'anomalies par arbres
iso = IsolationForest(contamination=0.05, random_state=42)
num_data = df[["Age", "Overall", "Potential"]].dropna()
labels = iso.fit_predict(num_data)

# -1 = outlier, 1 = inlier
n_outliers = (labels == -1).sum()
print(f"Isolation Forest outliers: {n_outliers} "
      f"({n_outliers / len(num_data) * 100:.1f}%)")
\end{minted}

% ============================================================
\section{R\`egles m\'etier}

\begin{minted}{python}
# Air Quality : la connaissance m\'etier compte
df_air = pd.read_csv("AirQualityUCI.csv", sep=";", decimal=",")

# Le dataset utilise -200 comme sentinelle pour donn\'ees manquantes
print(f"Values == -200: {(df_air == -200).sum().sum()}")

# Remplacer la sentinelle par NaN
df_air = df_air.replace(-200, np.nan)

# R\`egle m\'etier : une concentration de CO > 20 mg/m3
# est physiquement implausible en station urbaine
co_col = "CO(GT)"
implausible = df_air[co_col] > 20
print(f"Implausible CO readings: {implausible.sum()}")
df_air.loc[implausible, co_col] = np.nan
\end{minted}

\begin{datatip}
Les r\`egles m\'etier sont la m\'ethode de d\'etection d'outliers la plus fiable. Un z-score ne sait pas qu'un humain ne peut pas peser 800\,kg. Consultez toujours des experts m\'etier ou des plages de r\'ef\'erence quand elles existent.
\end{datatip}

% ============================================================
\section{Strat\'egies de traitement}

\begin{minted}{python}
# Strat\'egie 1 : supprimer les lignes aberrantes
df_clean = df[~outlier_mask].copy()

# Strat\'egie 2 : capper (winsorize) aux bornes
from scipy.stats import mstats
df["Overall_winsorized"] = mstats.winsorize(df["Overall"], limits=[0.01, 0.01])

# Strat\'egie 3 : remplacer par NaN et imputer plus tard
df.loc[outlier_mask, "Overall"] = np.nan

# Strat\'egie 4 : transformation log pour r\'eduire l'asym\'etrie
df["Value_log"] = np.log1p(df["Value(in Euro)"].clip(lower=0))

# Strat\'egie 5 : mise \`a l'\'echelle robuste (moins sensible aux outliers)
from sklearn.preprocessing import RobustScaler
scaler = RobustScaler()  # utilise m\'ediane et IQR au lieu de moyenne et std
df[["Overall_robust"]] = scaler.fit_transform(df[["Overall"]])
\end{minted}

\begin{minted}{python}
# Comparer les distributions avant et apr\`es winsorisation
fig, axes = plt.subplots(1, 2, figsize=(12, 4))

df["Overall"].hist(bins=50, ax=axes[0], alpha=0.7)
axes[0].set_title("Original")

df["Overall_winsorized"].hist(bins=50, ax=axes[1], alpha=0.7, color="green")
axes[1].set_title("Winsorized (1% each tail)")

plt.tight_layout()
plt.show()
\end{minted}

% ============================================================
\section{Exercices}

\begin{exercise}
\begin{enumerate}
  \item Chargez le dataset FIFA 21. Cr\'eez des box plots pour cinq colonnes num\'eriques. Identifiez par IQR les colonnes avec le plus d'outliers.
  \item Pour Air Quality UCI, remplacez toutes les sentinelles $-200$ par NaN. Puis appliquez des r\`egles m\'etier pour signaler les valeurs implausibles sur au moins deux colonnes de polluants.
  \item Comparez les nombres d'outliers par z-score et par IQR sur une colonne asym\'etrique (par ex. valeur de march\'e). Quelle m\'ethode est plus appropri\'ee et pourquoi ?
  \item Utilisez Isolation Forest sur Melbourne Housing (Price, Rooms, Distance, Landsize). Visualisez les outliers d\'etect\'es sur un scatter plot Price vs Landsize.
  \item Impl\'ementez une fonction \texttt{treat\_outliers(series, method)} qui supporte trois m\'ethodes : << remove >>, << cap >>, << log >>. Testez-la sur une colonne avec outliers connus.
\end{enumerate}
\end{exercise}

% ============================================================
\section{R\'esum\'e du chapitre}

\begin{itemize}
  \item Les outliers peuvent \^etre des erreurs, des artefacts ou de vraies valeurs extr\^emes --- l'investigation est obligatoire.
  \item Les m\'ethodes visuelles (box plots, scatter plots) donnent une intuition imm\'ediate.
  \item Le z-score marche pour des donn\'ees \`a peu pr\`es normales ; l'IQR est robuste aux distributions asym\'etriques.
  \item La distance de Mahalanobis et Isolation Forest d\'etectent des outliers multivari\'es invisibles aux m\'ethodes univari\'ees.
  \item Les r\`egles m\'etier (limites physiques, sentinelles) sont la m\'ethode la plus fiable.
  \item Les options de traitement : suppression, capping (winsorisation), transformation log, mise \`a l'\'echelle robuste.
\end{itemize}
